\chapter{Background and Related Work}\label{ch:background}
\thispagestyle{pagestyle}

\section{Reinforcement Learning Fundamentals}
\label{sec:rl-fundamentals}

Reinforcement learning is a computational framework for sequential
decision making under uncertainty \cite{Sutton2018}. An agent interacts
with an environment modelled as a Markov Decision Process (MDP), defined
by the tuple $(\mathcal{S}, \mathcal{A}, P, R, \gamma)$, where
$\mathcal{S}$ is the state space, $\mathcal{A}$ is the action space,
$P: \mathcal{S} \times \mathcal{A} \times \mathcal{S} \rightarrow [0,1]$
is the transition kernel, $R: \mathcal{S} \times \mathcal{A} \rightarrow
\mathbb{R}$ is the reward function, and $\gamma \in [0,1)$ is the
discount factor. The agent's goal is to learn a policy
$\pi: \mathcal{S} \rightarrow \mathcal{A}$ that maximises the expected
discounted return $G_t = \sum_{k=0}^{\infty} \gamma^k R_{t+k+1}$.

Policy gradient methods optimise the policy directly via the gradient of
the expected return with respect to policy parameters $\theta$.  The
fundamental result is the policy gradient theorem:
\begin{equation}
  \nabla_\theta J(\theta) = \mathbb{E}_{\pi_\theta}
  \left[ \nabla_\theta \log \pi_\theta(a|s) \cdot Q^{\pi_\theta}(s,a) \right].
  \label{eq:pg-theorem}
\end{equation}
Actor-critic architectures maintain both a policy network (actor) and a
value network (critic) trained simultaneously; a learned baseline
$V_\phi(s)$ reduces gradient variance through the advantage estimate
$A(s,a) = Q(s,a) - V_\phi(s)$.

\section{Proximal Policy Optimization}
\label{sec:ppo}

Proximal Policy Optimization \cite{Schulman2017} takes the largest
possible improvement step on a policy without destabilising training. It
replaces Trust Region Policy Optimization's constrained optimisation with
a clipped surrogate objective:
\begin{equation}
  L^{\mathrm{CLIP}}(\theta) = \mathbb{E}_t \left[
    \min\!\left( r_t(\theta)\,\hat{A}_t,\;
      \mathrm{clip}\!\left(r_t(\theta),\,1{-}\varepsilon,\,1{+}\varepsilon\right)\hat{A}_t
    \right)
  \right],
  \label{eq:ppo-clip}
\end{equation}
where $r_t(\theta) = \pi_\theta(a_t|s_t) / \pi_{\theta_{\mathrm{old}}}(a_t|s_t)$
is the probability ratio, $\hat{A}_t$ is the estimated advantage, and
$\varepsilon$ (typically $0.2$) is the clipping threshold. The full
objective adds a value-function regression term and an entropy bonus:
\begin{equation}
  L(\theta) = L^{\mathrm{CLIP}}(\theta) - c_1 L^{\mathrm{VF}}(\theta)
             + c_2 H[\pi_\theta(\cdot|s_t)].
  \label{eq:ppo-full}
\end{equation}
PPO is well suited to the vehicle control domain because its continuous
Gaussian policy maps naturally onto a steer/throttle/brake action space,
and its on-policy rollout structure integrates cleanly with a
wall-clock-rate simulation running in a separate process. The
Stable-Baselines3 library \cite{Raffin2021} provides the production
implementation used in this project, validated against reference
benchmarks.

\section{Simulation-to-Real Transfer in Autonomous Driving}
\label{sec:sim2real}

Training autonomous driving policies in simulation rather than on physical
hardware is standard practice, motivated by safety, cost, and the ability
to reset episodes deterministically \cite{Kiran2021}. The primary risk is
the \textit{sim-to-real gap}: discrepancies between simulated and real
sensor readings, physics, and environmental conditions cause policies that
perform well in simulation to degrade on deployment.

Prominent platforms include CARLA \cite{Dosovitskiy2017}, an
Unreal-Engine-based open-source simulator modelling urban driving at high
visual fidelity, and the Unity ML-Agents Toolkit \cite{Juliani2020}, which
provides a general-purpose RL training framework tightly integrated with
the Unity Editor. SHILATE uses raw Unity physics without ML-Agents for
two reasons: it avoids an opaque intermediate API layer, and the training
process communicates exclusively through MQTT --- the same protocol used
by the target SDV stack --- so no additional bridge is required for
deployment.

\section{Software-Defined Vehicles and the VSS Data Model}
\label{sec:sdv-vss}

The COVESA Vehicle Signal Specification (VSS) \cite{COVESA2023} is a
hierarchical, vendor-neutral data model for vehicle signals. VSS paths
follow a dotted-name convention, e.g.\ \texttt{Vehicle.Speed} (km/h),
\texttt{Vehicle.Powertrain.CombustionEngine.Speed} (RPM), and
\texttt{Vehicle.Chassis.SteeringWheel.Angle} (degrees).

The Eclipse Kuksa databroker \cite{EclipseKuksa2024} implements a
VSS-aligned gRPC server acting as the vehicle signal hub. Producers write
to VSS paths; consumers subscribe to changes. Eclipse Velocitas
\cite{EclipseVelocitas2024} builds on Kuksa with a vehicle-application
SDK and deployment scaffold. In SHILATE, the Velocitas application
monitors the running RL agent for speed and RPM limit violations.

\section{MQTT and Publish-Subscribe Messaging}
\label{sec:mqtt}

MQTT 3.1.1 \cite{OASIS2014} is a lightweight publish-subscribe protocol
designed for constrained IoT environments. A central broker decouples
producers and consumers: a publisher pushes a message to a topic string
and all subscribers to that topic receive it with no direct network
coupling between the two parties. SHILATE uses QoS~0 (at-most-once
delivery) throughout, which minimises latency at the cost of possible
message loss --- acceptable in a training loop where the next observation
replaces the current one within tens of milliseconds. The MQTT protocol
is also natively supported by the Eclipse Leda SDV image, which ships
Mosquitto as a core service.

\section{Related Work}
\label{sec:related}

A substantial body of work addresses RL for autonomous driving within
simulation \cite{Kiran2021}. OpenAI's \texttt{gym-car-racing}
\cite{Brockman2016} provides a top-down racetrack benchmark but offers no
path to real vehicle deployment. CARLA-based RL studies
\cite{Dosovitskiy2017} achieve high visual realism but operate as closed
systems with no integration to SDV middleware. The Unity ML-Agents toolkit
\cite{Juliani2020} enables RL training within Unity but couples the Python
trainer tightly to Unity's side-channel API, making it difficult to reuse
the same environment against a physical vehicle.

From the SDV side, the Eclipse Kuksa and Velocitas communities have
demonstrated containerised vehicle applications that read VSS signals and
enforce policies, but without an RL training pipeline attached. SHILATE
contributes the missing link: a training-to-deployment chain that is
agnostic to both the simulation engine and the SDV stack as long as both
speak MQTT and VSS.
